\documentclass[11pt]{article}

\usepackage[T1]{fontenc}
\usepackage{amsmath}
\usepackage{amssymb}
\usepackage[margin=1in]{geometry}
\usepackage{microtype}

\newcommand{\R}{\mathbb{R}}
\newcommand{\E}{\mathbb{E}}
\newcommand{\Var}{\operatorname{Var}}
\newcommand{\Cov}{\operatorname{Cov}}

\title{Coordinate Scaling for Raw PenroseDebruim Patches}
\author{}
\date{}

\begin{document}
\maketitle

\begin{abstract}
This note derives the exact part of the coordinate second moment that grows
with \(K\) in the raw De Bruijn construction.  The result is
\[
  V_{\mathrm{index}}(K)
  =\frac{25}{6}K(K+1)
\]
per Cartesian coordinate.  The population analysis assumes
\(\E[x]=\E[y]=0\) by symmetry and therefore identifies variance with raw
second moment; neither patch means nor finite-sample means are subtracted.
A seeded \(100{,}000\)-condition sweep for every \(K=0,\ldots,5\) calibrates
the bounded residual to the simple common value \(5.65\), giving
\[
  s_{xy}(K)
  =
  \sqrt{\frac{25}{6}K(K+1)+5.65}.
\]
The coefficient \(25/6\) is exact; \(5.65\) is empirical.  A sixth uniform
input rotates all five grid directions through one half of their angular
spacing.  This preserves the \(x/y\) second moment while making the wrapped
output angle uniform on \([-\pi,\pi)\).
\end{abstract}

\section{Raw construction}

Write
\[
  \gamma=(\gamma_0,\ldots,\gamma_4)\in[0,1)^5
\]
for the five phase values, and independently draw
\[
  \theta_{\mathrm{offset}}\sim\operatorname{Uniform}[0,1].
\]
Let \(K\) be a nonnegative integer and define the mapped angular offset
\[
  \theta_0
  =\left(\theta_{\mathrm{offset}}-\frac12\right)\frac{\pi}{5}
  \in\left[-\frac{\pi}{10},\frac{\pi}{10}\right].
\]
The five unit normals are
\[
  u_j =
  \begin{pmatrix}\cos\theta_j\\\sin\theta_j\end{pmatrix},
  \qquad
  \theta_j=\frac{2\pi j}{5}+\theta_0,
  \qquad j=0,\ldots,4.
\]
For phase vector \(\gamma\), family \(j\) contains lines
\[
  u_j^T x=n-\gamma_j,
  \qquad k\in\{-K,\ldots,K\}.
\]
For every unordered pair \(i<j\), and every ordered pair of line indices
\((k_i,k_j)\), one rhombus is emitted.  Hence
\[
  N(K)=\binom52(2K+1)^2
  =10(2K+1)^2.
\]
This includes \(N(0)=10\).

Let \(x=x_{ij}\) be the intersection of the selected lines.  The five region
coordinates are
\[
  r_i=k_i,\qquad r_j=k_j,\qquad
  r_m=\left\lceil u_m^T x+\gamma_m\right\rceil
  \quad(m\ne i,j).
\]
The raw rhombus center is
\begin{equation}
  c_{ij}
  =
  \sum_{m=0}^4 r_m u_m
  +\frac{u_i+u_j}{2}.
  \label{eq:raw-center}
\end{equation}
No patch mean is subtracted from \eqref{eq:raw-center}.  In particular, a
phase-dependent translation of the complete patch is part of the target.

\section{Exact line-index variance}

Let \(\mathbf{k}\) be uniform on the \(2K+1\) integers
\(\{-K,\ldots,K\}\).  Symmetry gives \(\E\mathbf{k}=0\), and
\[
  \Var(\mathbf{k})
  =\frac{1}{2K+1}
    \sum_{k=-K}^{K}k^2
  =\frac{K(K+1)}{3}.
  \label{eq:index-variance}
\]
This identity is exact, including at \(K=0\).

For a fixed family pair define
\[
  A_{ij}=
  \begin{pmatrix}u_i^T\\u_j^T\end{pmatrix},
  \qquad
  q_{ij}=
  \begin{pmatrix}k_i-\gamma_i\\k_j-\gamma_j\end{pmatrix}.
\]
Then
\[
  x_{ij}=A_{ij}^{-1}q_{ij}.
\]
Holding \(\gamma\) fixed and varying the independently enumerated indices,
the index-induced covariance is
\[
  \Cov_{\mathrm{indices}}(x_{ij})
  =\sigma_{\mathbf{k}}^2 A_{ij}^{-1}A_{ij}^{-T},
  \qquad
  \sigma_{\mathbf{k}}^2=\frac{K(K+1)}{3}.
\]
If \(\Delta_{ij}=\theta_j-\theta_i\), then
\[
  A_{ij}A_{ij}^T
  =
  \begin{pmatrix}
    1&\cos\Delta_{ij}\\
    \cos\Delta_{ij}&1
  \end{pmatrix}
\]
and therefore
\[
  \operatorname{tr}
  \left(A_{ij}^{-1}A_{ij}^{-T}\right)
  =2\csc^2\Delta_{ij}.
\]
After using rotational balance to divide the trace equally between \(x\) and
\(y\), the shared per-coordinate intersection contribution is
\(\sigma_{\mathbf{k}}^2\csc^2\Delta_{ij}\).

\section{Average over the ten family pairs}

Five unordered pairs have separation \(72^\circ\) modulo \(180^\circ\), and
five have separation \(144^\circ\), whose squared sine equals that of
\(36^\circ\).  Consequently
\begin{align*}
  \frac1{10}\sum_{i<j}\csc^2\Delta_{ij}
  &=
  \frac12\left(\csc^2 72^\circ+\csc^2 36^\circ\right)\\
  &=
  \frac12\left(
    \frac{8}{5+\sqrt5}+\frac{8}{5-\sqrt5}
  \right)
  =2,
\end{align*}
where
\[
  \sin^2 72^\circ=\frac{5+\sqrt5}{8},
  \qquad
  \sin^2 36^\circ=\frac{5-\sqrt5}{8}.
\]
Thus the pair-averaged intersection variance contributed by the line indices
is exactly
\[
  2\sigma_{\mathbf{k}}^2
  =\frac{2}{3}K(K+1)
\]
per coordinate.

\section{The dual-projection factor}

The five normals form a unit-norm tight frame:
\begin{equation}
  \sum_{m=0}^4u_m u_m^T=\frac52 I_2.
  \label{eq:tight-frame}
\end{equation}
This follows by expanding the matrix entries and observing that the sums of
\(\cos 2\theta_m\) and \(\sin 2\theta_m\) vanish.  Hence
\[
  \sum_{m=0}^4(u_m^Tx)u_m=\frac52x.
\]
The unquantized dual projection therefore multiplies distances by \(5/2\)
and coordinate variances by \(25/4\).  Its exact line-index contribution is
\begin{equation}
  V_{\mathrm{index}}(K)
  =
  \frac{25}{4}\,
  \frac{2}{3}K(K+1)
  =
  \boxed{\frac{25}{6}K(K+1)}.
  \label{eq:leading}
\end{equation}

\section{Exact decomposition of the bounded terms}

For each non-intersecting family define the ceiling remainder
\[
  \delta_m
  =
  \left\lceil u_m^Tx+\gamma_m\right\rceil
  -(u_m^Tx+\gamma_m),
  \qquad 0\le\delta_m<1
\]
up to the immaterial endpoint convention.  For the intersecting families,
the line equations give
\[
  r_i=u_i^Tx+\gamma_i,\qquad
  r_j=u_j^Tx+\gamma_j.
\]
Substituting these identities into \eqref{eq:raw-center} and applying
\eqref{eq:tight-frame} yields the exact decomposition
\begin{equation}
  c_{ij}
  =
  \frac52x_{ij}
  +\underbrace{\sum_{m=0}^4\gamma_m u_m}_{G(\gamma)}
  +\underbrace{\sum_{m\ne i,j}\delta_m u_m}_{Q_{ij}}
  +\frac{u_i+u_j}{2}.
  \label{eq:decomposition}
\end{equation}
The first term contains the exact growing contribution
\eqref{eq:leading}.  The remaining terms are bounded for fixed phases, but
bounded does not mean variance-free or independent:
\begin{itemize}
  \item \(G(\gamma)\) translates the raw patch as the phases change;
  \item each \(\delta_m\) is a sawtooth function of the intersection and
    is correlated with \(x_{ij}\);
  \item the set of omitted indices in \(Q_{ij}\) depends on the family pair;
  \item the half-diagonal \((u_i+u_j)/2\) is pair-dependent; and
  \item cross-covariances among all terms in \eqref{eq:decomposition}
    contribute to the total coordinate variance.
\end{itemize}

\section{Population moments and finite-sample estimators}

Let \(C=(X,Y)\) be the center obtained by first drawing
\(\Gamma\sim\operatorname{Uniform}([0,1)^5)\) and then selecting one tile
uniformly from its complete raw patch.  The population analysis adopts the
fivefold-symmetry identities
\[
  \E[X]=\E[Y]=0.
\]
It follows that
\[
  \Var(X)=\E[X^2],\qquad
  \Var(Y)=\E[Y^2],
\]
and the shared population scale is defined from the pooled raw second moment:
\begin{equation}
  s_{xy}^2(K)
  =
  \frac{\E[X^2]+\E[Y^2]}{2}.
  \label{eq:pooled-second-moment}
\end{equation}
Neither a generated patch mean nor a finite Monte Carlo sample mean is
subtracted in \eqref{eq:pooled-second-moment}.

For auditability, the sweep also records centered empirical quantities.  With
\(M\) observed tiles,
\[
  \bar x=\frac1M\sum_{m=1}^M x_m,\qquad
  v_x=\frac1M\sum_{m=1}^M(x_m-\bar x)^2,\qquad
  q_x=\frac1M\sum_{m=1}^M x_m^2.
\]
The columns \texttt{variance\_x} and \texttt{sample\_variance\_x} use
denominators \(M\) and \(M-1\), respectively; the same definitions apply to
\(y\) and wrapped angle.  The pooled \(xy\) sample is the concatenation
\((x_1,\ldots,x_M,y_1,\ldots,y_M)\).  Therefore
\[
  \overline{xy}
  =\frac{\sum_m x_m+\sum_m y_m}{2M},
  \qquad
  q_{xy}
  =\frac{\sum_m x_m^2+\sum_m y_m^2}{2M}.
\]
Its centered population and sample variances use denominators \(2M\) and
\(2M-1\).  The pooled mean is retained as a diagnostic even though separate
\(\bar x\) and \(\bar y\) are more geometrically informative.

The observed \(x\)- and \(y\)-means are at most about \(7.0\times10^{-4}\) in
absolute value.  Under the zero-population-mean symmetry assumption these are
finite-sample estimation errors.  They are reported but are not used to
center \(q_x\), \(q_y\), or \(q_{xy}\).

\section{Raw second-moment sweep and residual fit}

The reproducible sweep uses seed \(20260823\), \(100{,}000\) independent
independent \((\gamma,\theta_{\mathrm{offset}})\) draws for each
\(K=0,\ldots,5\), and every tile in every raw
patch.
Thus the tile observation count is
\[
  M=100{,}000\cdot10(2K+1)^2.
\]
The full statistics are stored in
\texttt{tests/scaling/raw\_statistics.csv}.  The entries relevant to the
coordinate scale are:
\begin{center}
\begin{tabular}{rrrr}
\(K\)&\(q_{xy}\)&exact leading term&residual\\
\hline
0&5.653308993&0&5.653308993\\
1&13.980624932&8.333333333&5.647291599\\
2&30.645533604&25&5.645533604\\
3&55.646613015&50&5.646613015\\
4&88.978783742&83.333333333&5.645450409\\
5&130.646957465&125&5.646957465
\end{tabular}
\end{center}
Here ``residual'' means
\[
  R_K^{\mathrm{emp}}
  =
  q_{xy}-
  \frac{25}{6}K(K+1).
\]
The exact analysis does not prove that this residual is constant.  Empirically,
however, all six values lie close to \(5.65\).  The simplest calibrated model
over the requested domain is a constant residual.  Its least-squares estimate
in variance space is the arithmetic mean \(5.647525847\); rounding to three
significant digits gives
\[
  R_{\mathrm{cal}}=5.65.
\]
This is deliberately not presented with spurious extra digits.  Combining it
with \eqref{eq:leading} gives the operational formula
\begin{equation}
  \boxed{
  s_{xy}(K)
  =
  \sqrt{\frac{25}{6}K(K+1)+5.65}
  }.
  \label{eq:scale}
\end{equation}
Its largest relative discrepancy from the six measured second-moment scales
is approximately \(0.03\%\), at \(K=0\); for
\(K=1,\ldots,5\) it is below \(0.01\%\).  Runtime code evaluates
\eqref{eq:scale} directly and never reads the statistics CSV.

\section{Wrapped angle statistics}

The CSV angle columns refer to the raw generated angle wrapped into
\([-\pi,\pi)\), in radians, before the training multiplier
\(\sqrt3/\pi\) is applied.

Under the requested equal-count assumption, the construction produces ten
equally weighted directions separated by \(\pi/5\).  The new independent
variable rotates all of those directions by
\[
  \theta_0\sim\operatorname{Uniform}
  \left[-\frac{\pi}{10},\frac{\pi}{10}\right].
\]
The intervals swept by adjacent directions meet exactly at their endpoints.
Their equally weighted union is therefore the uniform distribution on the
whole wrapped circle:
\[
  \Theta\sim\operatorname{Uniform}[-\pi,\pi).
\]
Consequently,
\[
  \E[\Theta]=0,\qquad
  \Var(\Theta)=\E[\Theta^2]=\frac{\pi^2}{3}.
\]
After training scaling, the second moment is multiplied by
\[
  \left(\frac{\sqrt3}{\pi}\right)^2=\frac3{\pi^2}.
\]
Thus scaled angle has exact unit second moment in the population.  The
committed angle-histogram CSV divides
\([-\pi,\pi)\) into twenty intervals of width \(\pi/10\); each interval has
population probability \(1/20\).  Finite-sample counts and means fluctuate
around \(5\%\) and zero, respectively.

The global rotation also preserves \(\|c_{ij}\|^2=x^2+y^2\) for every tile.
Therefore it does not change the pooled \(xy\) second moment or the scale
derived in the preceding sections.

\section{Status summary}

\begin{itemize}
  \item \(N(K)=10(2K+1)^2\) is exact.
  \item The generator emits raw centers and never subtracts a patch mean.
  \item Population \(x/y\) scaling uses raw second moments under
    \(\E[X]=\E[Y]=0\).
  \item The growing term \((25/6)K(K+1)\) is exact.
  \item The common residual \(5.65\) is calibrated from the six seeded raw
    sweeps and is not an analytic identity.
  \item Runtime scale is \eqref{eq:scale}; the CSV is statistics only.
  \item The sixth uniform input rotates all directions by
    \([-\pi/10,\pi/10]\), making wrapped angle uniform with zero mean and
    second moment \(\pi^2/3\).
\end{itemize}

\end{document}
